\documentclass[../main.tex]{subfiles}
\begin{document}

% Nguồn: docs/PROGRESS.md (bảng trạng thái Task 1--15, rulings), PLAN.md (Phase 5--7, mục 8),
% docs/superpowers/specs/2026-10-25-proxy-benchmark-loop-design.md, config/site.yaml, config/llm.yaml;
% cập nhật 2026-10-31: report/2026-10-27/data/campaign_specs.csv, experiments/campaigns/REF_run.log, report/2026-10-28, report/2026-10-29, REF_n29s5;
% cập nhật 2026-11-09: report/2026-11-08/README.md (REF_n1104s5, B2/P1/P2_n1108, AB_cut, AB_d1, AB_d1v).

Chương này tổng kết những gì luận văn đã làm, những gì quá trình xây dựng cho thấy, những việc chưa làm, và hướng phát triển tiếp theo.

\section{Kết luận}

\subsection{Mục tiêu và cách làm}

Luận văn mô phỏng hành vi tham gia giao thông tại nút Lê Văn Lương -- Hoàng Đạo Thúy trong cao điểm sáng, sao cho giống thực tế nhất có thể. Kịch bản gồm 11 hành vi; rẽ trái trái phép chờ giữa nút là một trong số đó. Một LLM chạy cục bộ sinh và tinh chỉnh lặp các luật hành vi tài xế cho SUMO~\cite{lopez2018sumo}. LLM chỉ sinh đặc tả luật ngoài vòng mô phỏng. SUMO thực thi phần tham số vType tĩnh, và một bộ điều khiển TraCI thực thi phần luật có điều kiện.

\subsection{Đã hoàn thành}

Toàn bộ mã của quy trình đã chạy được và có 403 ca kiểm thử tự động; trên máy chính ngày 08/11/2026, 401 ca đạt và 2 ca xfail có chủ đích (Mục~\ref{section:testing}).

Mạng mô phỏng dựng hoàn toàn bằng script từ OpenStreetMap qua tệp vá sinh từ cấu hình hiện trường, với đèn ba pha, làn rẽ phải liên tục ở 30~m cuối của mọi nhánh, và các bản vá chỉ cho lớp xe vi phạm (kết nối rẽ trái trái phép ở hai nhánh Hoàng Đạo Thúy, hai làn vỉa hè). Cấu hình cơ sở B0, bộ sinh nhu cầu, trình chạy kèm \texttt{manifest.json} và bộ chuyển đầu ra sang lược đồ quan sát chung chạy được trên cả Windows và Linux; tám lần chạy B0 đối chiếu giữa hai máy cho cùng số va chạm và teleport.

Phần luật gồm lược đồ đặc tả, tệp miền giá trị có ghi nguồn cho 24 thuộc tính vType, bộ kiểm tra hợp lệ trả lỗi có cấu trúc theo tám mã, bộ biên dịch, và bộ điều khiển TraCI với sáu luật có điều kiện. Trong đó \texttt{queue\_space\_seeking} gộp các hành vi tìm khoảng trống và có lựa chọn đi sai làn kiểu cắt về; luật vượt đỏ chỉ cho xe vào nút khi nút thông. Phần đánh giá gồm chỉ số M1--M10, cổng GEH, hàm khoảng cách có chặn trên và điểm tổng $S$, cùng khung đánh giá dùng chung với bộ đếm ngân sách theo chiến dịch. Phần hiệu chỉnh gồm trình khách LLM với năm phiên bản prompt, vòng lặp LLM chạy tiếp được sau khi bị ngắt, và phương án Optuna. Các đóng góp kỹ thuật rút ra từ quá trình này nằm ở Chương~\ref{chapter:SolutionAndContribution}.

\subsection{Phát hiện trong quá trình xây dựng}

Nhiều hiện tượng trông giống hành vi tài xế hóa ra đến từ cách dựng mạng. Va chạm cao bất thường của B0 trên mạng ban đầu đến từ bề rộng mặc định không có nguồn của làn ra. Trên mạng 29/10, va chạm trên làn của B0 là 18,53/h; trên mạng 04/11, chỉ bỏ các kết nối chỉ cho xe vi phạm từ lane 0, 1 đã làm mức này còn 1,87/h và B0 đủ điều kiện ở 15/15 lần chạy, dù không xe nào của B0 đi các kết nối đó. Cơ chế của ảnh hưởng này chưa xác định.

Bật kiểm tra va chạm trong nút phát hiện một phiên bản luật giữ xe tưởng an toàn nhưng thực chất cho xe đi xuyên qua nhau. Từ đó có nguyên tắc luật TraCI chỉ giữ xe, không đẩy xe qua xung đột. Theo cùng tinh thần, luật vượt đỏ được thêm điều kiện nút thông sau khi chẩn đoán cho thấy nó là điều kiện cần cho các lần khóa nút của B1 trên mạng 29/10; teleport kẹt của B1 giảm từ 51,73/h xuống 4,67/h trên 15 lần chạy.

Các hành vi tìm khoảng trống phụ thuộc mạnh vào cách mô hình hóa không gian. SUMO không tự tạo hành vi đi trên vỉa hè, nhưng biểu diễn được bằng cấu trúc mạng riêng kèm điều kiện kích hoạt qua TraCI. Hành vi đi ngược chiều chưa biểu diễn được mà vẫn giữ việc tính xung đột với dòng ngược chiều, nên được hoãn. Khi lane 0 cho đi thẳng theo đèn như ngoài hiện trường, hành vi chặn làn rẽ phải xuất hiện ngay ở B0 (8,18 xe đứng chờ đỏ mỗi chu kỳ, so với 0,57 của dữ liệu đại diện); mục tiêu M9 chưa được định nghĩa lại theo mô hình này. Xe đi sai làn kiểu cắt về đã nhập lại làn đúng trước vạch dừng, nên M8 phải đọc làn ở đầu đoạn 30~m thay vì tại vạch.

Hai phát hiện liên quan trực tiếp tới các chiến dịch. Thứ nhất, xe chờ chèn ở PROXY\_D1 của mọi đặc tả dựa trên cấu trúc B1 chủ yếu đến từ \texttt{accel} 1,1~m/s\textsuperscript{2} của xe máy, chuyển từ tham số IDM của Saval và cộng sự~\cite{saval2023mixed}; bỏ hết luật TraCI vẫn còn 203--387 xe chờ chèn, còn trả \texttt{accel} về mặc định SUMO thì còn 10--126 xe nhưng va chạm tăng lên 48--58/h. Thứ hai, \texttt{qwen3.5:9b} mắc chủ yếu lỗi cấu trúc (mã nguyên mẫu trùng tên phân phối, tổng tỷ trọng), và khi đã có một cấu trúc hợp lệ trong lịch sử thì giữ nguyên nó; ở chế độ sinh cả cấu trúc, mô hình hạ $S$ bằng cách bỏ bớt luật.

\subsection{Kết quả tạm thời của các chiến dịch}

Trên mạng F2, Optuna B2 đã dùng hết 36 lần chạy; không đặc tả nào trong 12 đặc tả của nó đủ điều kiện vì tất cả còn teleport do kẹt, và $S$ của chúng nằm trong 1,182--1,523, so với 1,181 của B0 chấm cùng cách (Mục~\ref{subsection:interim}). Trên hai mạng ngày 28/10, B2 dùng 24 lần chạy mỗi mạng; không đặc tả nào đủ điều kiện, và 3 trong 16 đặc tả có $S$ nhỏ hơn B0 của cùng mạng (Mục~\ref{subsection:scenario_update}).

Trên mạng 04/11, B0 và B1 chạy 15 lần mỗi cấu hình: B0 đủ điều kiện ở 15/15 lần với $S = 1{,}2174$, còn B1 bị loại ở cả 15 lần với $S = 1{,}2494$ (Mục~\ref{subsection:ref_n1104}). Ba chiến dịch B2, P1 và P2 (P1, P2 với \texttt{qwen3.5:9b}) mỗi chiến dịch dùng 36 lần chạy trên PROXY\_D1--D3 với seed 1, và không có đặc tả nào đủ điều kiện. $S$ nhỏ nhất là 1,1438 với B2, 1,2386 với P1 và 1,1680 với P2. P1 không thể đủ điều kiện vì giữ cấu trúc B1; P2 có 11/24 lần chạy đủ điều kiện ở PROXY\_D2 và PROXY\_D3, nhưng các đặc tả của nó bỏ \texttt{queue\_space\_seeking}, \texttt{uturn\_bypass} và \texttt{slow\_lane\_center} (Mục~\ref{subsection:campaigns_n1108}).

Các số này dựa trên một seed mỗi ngày, trên dữ liệu đại diện \texttt{LITERATURE\_PROXY} và các giá trị tạm thời. Chúng chưa cho phép kết luận LLM tốt hơn hay kém hơn Optuna.

\subsection{Chưa thực hiện}

Luận văn chưa có quyết định về \texttt{accel} của xe máy trong B1, chưa chạy lại P2 với prompt \texttt{v005}, chưa chạy lại các đặc tả chọn được với 5 seed, chưa chạy phương án H, và chưa có đánh giá cuối E1, E2 trên PROXY\_D4--PROXY\_D5 với ít nhất 20 seed và kiểm định thống kê. Luận văn cũng chưa có dữ liệu quan sát từ video; mọi giá trị đo trong cấu hình hiện trường thực vẫn để trống, và mọi giá trị tạm thời in màu đỏ trong luận văn còn chờ thay. Vì vậy luận văn chưa khẳng định mô phỏng tái hiện đúng giao thông thực tại nút. \textbf{[TODO: kết quả P2 với \texttt{v005}, kết quả H]} \textbf{[TODO: kết quả đánh giá cuối E1, E2 trên PROXY\_D4--PROXY\_D5 với kiểm định thống kê]} \textbf{[TODO: câu kết luận chính về khả năng của LLM cục bộ trong hiệu chỉnh luật hành vi, chỉ trong phạm vi dữ liệu đại diện]}

\section{Hướng phát triển}

\subsection{Việc trước mắt}

Việc đầu tiên là quyết định giữ \texttt{accel} xe máy của Saval trong B1 hay sửa B1, bằng cách dùng mặc định SUMO cho thuộc tính này hoặc nâng cận dưới của nó trong tệp miền giá trị (Mục~\ref{subsection:d1_backlog}). Khi PROXY\_D1 còn xe chờ chèn, không đặc tả nào dựa trên cấu trúc B1 đủ điều kiện được, nên mọi chiến dịch sau phụ thuộc vào quyết định này. Tiếp theo, P2 được chạy lại với prompt \texttt{v005} để đặc tả có đủ sáu luật TraCI, và các đặc tả chọn được của B2 và P2 được chạy lại với 5 seed trước khi so với B0 và B1. Sau đó, phương án H chạy Optuna trên cấu trúc tốt nhất của P2 với cùng ngân sách 36 lần chạy. Cuối cùng, B0, B1, B2, P1, P2 và H được đánh giá trên PROXY\_D4--PROXY\_D5 với ít nhất 20 seed, dùng kiểm định Mann--Whitney có hiệu chỉnh Holm và kích thước hiệu ứng Cliff's delta cho E1 (P1 so với B2) và E2 (B2, H, P2).

Song song, vòng LLM có hai cải tiến đã xác định: dùng lại kết quả khi LLM nộp lại một đặc tả đã chạy (ở P1, ba lần nộp lại tốn 9 lần chạy), và ép trong lược đồ tỷ trọng bằng 1 cho phân phối chỉ có một nguyên mẫu. Hai vấn đề phương pháp cần quyết định trước khi diễn giải kết quả là cách xử lý tỷ lệ nhỏ trong hàm khoảng cách, để bộ tối ưu không lợi khi ép một hành vi hiếm về 0, và định nghĩa lại hai mục tiêu của M9 theo mô hình làn 30~m.

\subsection{Dữ liệu đo thực}

Hướng quan trọng nhất về mặt khoa học là thay dữ liệu đại diện bằng dữ liệu đo thực, tức thay toàn bộ các giá trị in màu đỏ. Việc đầu tiên là quay video tại nút trong nhiều ngày và xây một đường ống xử lý video (ngoài phạm vi luận văn) xuất đúng lược đồ quan sát đã định nghĩa (Phụ lục~\ref{section:appendix_schema}), kể cả các bảng của nhóm tìm khoảng trống và làn của xe tại đầu đoạn làn rẽ phải cho M8. Khi đó các ngày PROXY được thay bằng ngày đo, việc tách ngày hiệu chỉnh và ngày đánh giá độc lập thực sự, và sàn nhiễu giữa các ngày tính từ biến động quan sát thay vì hệ số giả định.

Việc thứ hai là đo hoặc xác định tại hiện trường các giá trị đang để trống trong cấu hình hiện trường thực: bề rộng từng làn vào và làn ra; hướng đi được phép của từng làn; các nhánh cấm rẽ trái hoặc quay đầu và các làn xe vi phạm dùng để rẽ; chu kỳ và các pha đèn, gồm vàng và đỏ toàn phần; chiều dài đoạn làn rẽ phải liên tục; vị trí và bề rộng vỉa hè cùng điểm xe lên vỉa hè và điểm nhập lại. Việc thứ ba là chạy lại, trên hình học đo thực, thí nghiệm sức chứa vùng chờ giữa nút ở hai nhánh Hoàng Đạo Thúy và thí nghiệm 3B, và thay phạm vi quan sát hàng chờ hiện đặt 80~m bằng phạm vi thực của camera. Mọi thực nghiệm của Chương~\ref{chapter:Experiment} sau đó được chạy lại trên các giá trị thật.

\subsection{Mở rộng}

Khi bộ luật đã hiệu chỉnh trên dữ liệu thực, mô phỏng có thể dùng cho các kịch bản chính sách, ví dụ đổi lưu lượng rẽ trái trái phép ở 0\%, 50\% và 100\% so với quan sát để đo độ trễ và số xung đột có PET dưới ngưỡng. Kết quả loại này phải ghi rõ là ngoại suy. Các mở rộng khác gồm dùng thiết bị đo xung đột SSM của SUMO cho chẩn đoán an toàn trên các lần chạy đánh giá; kiểm tra luật cắt về trên dữ liệu video, so tỷ lệ xe đi sai làn và vị trí xe nhập lại với quan sát; tìm cách biểu diễn hành vi đi ngược chiều mà vẫn giữ việc tính xung đột với dòng ngược chiều thật; và mở rộng sang các nút giao khác để kiểm tra khả năng chuyển giao của quy trình.

\end{document}
